Subjects i =1...n
Test items j= 1...m

each response y_ij  =1 if correct else 0

there is some latent trait \xi_i, usually in R^1  

PCA
IRT has more theory, and is connected to the binary setting

teaches us about subjects and item parameters.

ordinal data

p_ij\equiv Pr(y_ij=1)=F(\xi_i\beta_i-a_j)

\beta is the item discrimination paramter
a_j is the item difficulty parameter.

F: is like a logit or something.

a_j/\beta_j is cutpoint, which is the .5 level, and you are equally likely to get question right or wrong.  these are the questions that they give you.
lots of parameters:
n*d ability parameters
m*(d+1)


\xi_i\beta_j is unidentified, set \xi to have mean 0, variance 1.    Scale relative to a reference item, or set a discrimination parameter on some of the questions.

NOMINATE
"alternating maximum likelihood"
estimate conditional on item parameters
estimate item parameters on abilities.

Incidental parameter problem

EM 
vs 
MCMC
M step Regress Y on X
E step E(Y|x\beta y-1)
Bayes visits the location proportional to the posterior